\chapter{Tremors}

\section*{Definition}
Tremors are involuntary, rhythmic, oscillatory movements of a body part. They are the most common movement disorder. Tremors are classified by their appearance (resting, action, intention) and underlying cause.

\section*{What Happens in Your Body}
Tremors arise from dysfunction in the central nervous system circuits that control movement. Resting tremor (Parkinson disease) involves degeneration of dopaminergic neurons in the substantia nigra. Essential tremor involves abnormal oscillations in the cerebello-thalamo-cortical circuit. Cerebellar tremor results from damage to the cerebellum or its connections. Enhanced physiologic tremor occurs when the normal tremor is amplified by metabolic or drug effects.

\section*{Causes}
\begin{itemize}
\item Essential tremor (most common movement disorder, often familial)
\item Parkinson disease (resting tremor, bradykinesia, rigidity)
\item Enhanced physiologic tremor: anxiety, stress, fatigue, caffeine, stimulants
\item Medication-induced: beta-agonists, antidepressants, valproate, lithium, corticosteroids
\item Metabolic: hyperthyroidism, hypoglycemia, electrolyte imbalances
\item Alcohol withdrawal
\item Cerebellar tremor: multiple sclerosis, stroke, tumor, trauma
\item Drug-induced: alcohol, benzodiazepine withdrawal
\item Wilson disease (young-onset, with other neurological signs)
\end{itemize}

\section*{When to See a Doctor}
\begin{itemize}
\item Rhythmic, involuntary shaking of a body part
\item Resting tremor (present at rest, decreases with movement - Parkinson)
\item Action tremor (occurs with voluntary movement - essential tremor)
\item Intention tremor (worsens at the end of a goal-directed movement - cerebellar)
\item Postural tremor (occurs when maintaining a position against gravity)
\item Tremor that interferes with writing, eating, drinking, or other daily tasks
\item Anxiety or social embarrassment due to visible tremor
\end{itemize}

\section*{Related Symptoms}
\begin{itemize}
\item Parkinson disease
\item Essential tremor
\item Multiple sclerosis
\item Hyperthyroidism
\item Anxiety
\item Alcohol withdrawal
\end{itemize}

\section*{Self-Care / Home Management}
\begin{itemize}
\item Reduce or avoid caffeine, nicotine, and other stimulants
\item Practice relaxation techniques to reduce stress
\item Get adequate sleep
\item Use weighted utensils, cups, and pens to dampen tremor
\item Avoid alcohol (temporary improvement can lead to dependence)
\item Use adaptive devices (tremor-canceling spoons, non-slip mats)
\end{itemize}

\section*{How Your Doctor Will Evaluate You}
\begin{itemize}
\item Clinical history and physical examination
\item Differentiate rest vs. action vs. intention tremor
\item Blood tests: thyroid function, glucose, electrolytes, liver function
\item Brain MRI to rule out structural lesions
\item DaTscan (dopamine transporter scan) for Parkinson disease
\item Genetic testing for familial essential tremor
\end{itemize}

\section*{Treatment Options}
\begin{itemize}
\item Essential tremor: beta-blockers (propranolol), primidone; deep brain stimulation for refractory
\item Parkinson tremor: levodopa, dopamine agonists; deep brain stimulation
\item Enhanced physiologic tremor: treat underlying cause
\item Cerebellar tremor: physical therapy, weighted limbs; limited medication options
\item Botulinum toxin injections for head or voice tremor
\item Deep brain stimulation for severe, medication-refractory tremor
\end{itemize}

\section*{References}
\begin{thebibliography}{99}
\bibitem{tremors:uptodate-tremor} Elble RJ, et al. Tremor: clinical features and diagnosis. In: UpToDate, Post TW (Ed), UpToDate, Waltham, MA; 2025.
\bibitem{tremors:nejm-essential-tremor} Louis ED, et al. ``Essential Tremor.'' \emph{N Engl J Med}. 2023;389(13):1223-1233.
\bibitem{tremors:lancet-neurol-tremor} Bhatia KP, et al. ``Tremor: clinical classification.'' \emph{Lancet Neurol}. 2024;23(3):278-290.
\bibitem{tremors:bmj-tremor} Deuschl G, et al. ``Tremor in adults.'' \emph{BMJ}. 2024;386:e082143.
\bibitem{tremors:cochrane-tremor} Zesiewicz TA, et al. ``Pharmacological treatment of essential tremor.'' \emph{Cochrane Database Syst Rev}. 2023;(10):CD013576.
\bibitem{tremors:harrison-tremor} Longo DL, et al. \emph{Harrison\'s Principles of Internal Medicine}. 21st ed. McGraw-Hill; 2022. Chapter 418: Tremor.
\end{thebibliography}
